\section{Architecture Taxonomy and Implementation}
\label{sec:architecture-taxonomy}

\subsection{Attention and Sequence-Mixer Families}
\label{sec:attention-families}

This system implements thirty-five sequence mixer variants organized into six functional categories reflecting research trends in sequence modeling design. The taxonomical organization reflects evolving understanding of how to balance expressivity, computational efficiency, and memory constraints.

\begin{enumerate}[leftmargin=*]
    \item \textbf{Dense Attention Baselines} (3): Standard softmax attention, sigmoid attention, and Grouped-Query Attention (GQA) provide full global contextualization at quadratic computational cost, serving as reference baselines for comparison with more efficient alternatives.
    \item \textbf{Recurrent and Retentive Architectures} (5): RetNet, Mamba, ODE-style blocks, Titans, and Engram maintain state representations enabling $\mathcal{O}(1)$ inference cost while preserving expressivity through recurrent dynamics, selective parameters, or test-time memory adaptation.
    \item \textbf{Sparse Attention Patterns} (9): Sparse Transformer, Longformer, BigBird, SparseK, NSA, SpargeAttn, FASA, MSA, and SparDA reduce quadratic complexity through structured sparsity, token selection, or training-free pruning strategies.
    \item \textbf{Gated Memory Mechanisms} (8): GLA, DeltaNet, Gated DeltaNet, Gated DeltaNet-2, HGRN2, FoX, Gated Softmax, and KDA introduce data-dependent control over memory retention, forgetting, and update strength.
    \item \textbf{Latent and Low-Rank Compression} (9): MLA, GQLA, MLRA, Tucker, IHA, GTA, MTLA, CCA, and CCGQA compress key--value representations into low-rank latent bottlenecks, trading a small reconstruction error for substantial memory and compute savings.
    \item \textbf{Conditional Memory} (1): Engram provides test-time memory through learned read/write operations over an external memory bank.
\end{enumerate}

\begin{figure}[t]
\centering
\fitfigure{\begin{tikzpicture}[
    font=\scriptsize,
    node distance=0.32cm and 0.55cm,
    box/.append style={inner sep=3pt, minimum height=0.48cm}
]
% Root
\node[box, fill=blue!15, minimum width=4.0cm] (root) {\textbf{Sequence Mixer Registry (35+ variants)}};

% Category headers — 6 columns
\node[box, fill=red!10, below left=1.0cm and 7.5cm of root, minimum width=2.0cm] (dense) {\textbf{Dense (3)}};
\node[box, fill=green!10, below left=1.0cm and 4.5cm of root, minimum width=2.0cm] (recur) {\textbf{Recurrent (5)}};
\node[box, fill=yellow!12, below left=1.0cm and 1.5cm of root, minimum width=2.0cm] (sparse) {\textbf{Sparse (9)}};
\node[box, fill=purple!12, below right=1.0cm and 1.5cm of root, minimum width=2.0cm] (gated) {\textbf{Gated (8)}};
\node[box, fill=cyan!12, below right=1.0cm and 4.5cm of root, minimum width=2.0cm] (latent) {\textbf{Latent (9)}};
\node[box, fill=orange!12, below right=1.0cm and 7.5cm of root, minimum width=2.0cm] (memory) {\textbf{Memory (1)}};

% Dense column (3)
\node[box, fill=red!5, below=of dense] (d1) {\texttt{standard\_attn}};
\node[box, fill=red!5, below=of d1] (d2) {\texttt{sigmoid\_attn}};
\node[box, fill=red!5, below=of d2] (d3) {\texttt{gqa}};

% Recurrent column (5)
\node[box, fill=green!5, below=of recur] (r1) {\texttt{retnet / retnet\_attn}};
\node[box, fill=green!5, below=of r1] (r2) {\texttt{mamba}};
\node[box, fill=green!5, below=of r2] (r3) {\texttt{ode}};
\node[box, fill=green!5, below=of r3] (r4) {\texttt{titan\_attn}};
\node[box, fill=green!5, below=of r4] (r5) {\texttt{engram\_attn}};

% Sparse column (9)
\node[box, fill=yellow!6, below=of sparse] (s1) {\texttt{sparse\_transformer\_attn}};
\node[box, fill=yellow!6, below=of s1] (s2) {\texttt{longformer\_attn}};
\node[box, fill=yellow!6, below=of s2] (s3) {\texttt{bigbird\_attn}};
\node[box, fill=yellow!6, below=of s3] (s4) {\texttt{sparsek\_attn}};
\node[box, fill=yellow!6, below=of s4] (s5) {\texttt{nsa\_attn}};
\node[box, fill=yellow!6, below=of s5] (s6) {\texttt{sparge\_attn}};
\node[box, fill=yellow!6, below=of s6] (s7) {\texttt{fasa\_attn}};
\node[box, fill=yellow!6, below=of s7] (s8) {\texttt{msa\_attn}};
\node[box, fill=yellow!6, below=of s8] (s9) {\texttt{sparda\_attn}};

% Gated column (8)
\node[box, fill=purple!6, below=of gated] (g1) {\texttt{gla\_attn}};
\node[box, fill=purple!6, below=of g1] (g2) {\texttt{deltanet\_attn}};
\node[box, fill=purple!6, below=of g2] (g3) {\texttt{gated\_deltanet\_attn}};
\node[box, fill=purple!6, below=of g3] (g4) {\texttt{gated\_deltanet2\_attn}};
\node[box, fill=purple!6, below=of g4] (g5) {\texttt{hgrn2\_attn}};
\node[box, fill=purple!6, below=of g5] (g6) {\texttt{fox\_attn}};
\node[box, fill=purple!6, below=of g6] (g7) {\texttt{gated\_softmax\_attn}};
\node[box, fill=purple!6, below=of g7] (g8) {\texttt{kda\_attn}};

% Latent column (9)
\node[box, fill=cyan!6, below=of latent] (l1) {\texttt{mla\_attn}};
\node[box, fill=cyan!6, below=of l1] (l2) {\texttt{gqla\_attn}};
\node[box, fill=cyan!6, below=of l2] (l3) {\texttt{mlra\_attn}};
\node[box, fill=cyan!6, below=of l3] (l4) {\texttt{tucker\_attn}};
\node[box, fill=cyan!6, below=of l4] (l5) {\texttt{iha\_attn}};
\node[box, fill=cyan!6, below=of l5] (l6) {\texttt{gta\_attn}};
\node[box, fill=cyan!6, below=of l6] (l7) {\texttt{mtla\_attn}};
\node[box, fill=cyan!6, below=of l7] (l8) {\texttt{cca\_attn}};
\node[box, fill=cyan!6, below=of l8] (l9) {\texttt{ccgqa\_attn}};

% Memory column (1)
\node[box, fill=orange!6, below=of memory] (m1) {\texttt{engram\_attn}};

% Root-to-category links
\draw[line] (root) -- (dense);
\draw[line] (root) -- (recur);
\draw[line] (root) -- (sparse);
\draw[line] (root) -- (gated);
\draw[line] (root) -- (latent);
\draw[line] (root) -- (memory);

% Category-to-list links — Dense
\draw[line] (dense) -- (d1);
\draw[line] (d1) -- (d2);
\draw[line] (d2) -- (d3);

% Category-to-list links — Recurrent
\draw[line] (recur) -- (r1);
\draw[line] (r1) -- (r2);
\draw[line] (r2) -- (r3);
\draw[line] (r3) -- (r4);
\draw[line] (r4) -- (r5);

% Category-to-list links — Sparse
\draw[line] (sparse) -- (s1);
\draw[line] (s1) -- (s2);
\draw[line] (s2) -- (s3);
\draw[line] (s3) -- (s4);
\draw[line] (s4) -- (s5);
\draw[line] (s5) -- (s6);
\draw[line] (s6) -- (s7);
\draw[line] (s7) -- (s8);
\draw[line] (s8) -- (s9);

% Category-to-list links — Gated
\draw[line] (gated) -- (g1);
\draw[line] (g1) -- (g2);
\draw[line] (g2) -- (g3);
\draw[line] (g3) -- (g4);
\draw[line] (g4) -- (g5);
\draw[line] (g5) -- (g6);
\draw[line] (g6) -- (g7);
\draw[line] (g7) -- (g8);

% Category-to-list links — Latent
\draw[line] (latent) -- (l1);
\draw[line] (l1) -- (l2);
\draw[line] (l2) -- (l3);
\draw[line] (l3) -- (l4);
\draw[line] (l4) -- (l5);
\draw[line] (l5) -- (l6);
\draw[line] (l6) -- (l7);
\draw[line] (l7) -- (l8);
\draw[line] (l8) -- (l9);

% Category-to-list links — Memory
\draw[line] (memory) -- (m1);
\end{tikzpicture}}
\caption{Comprehensive taxonomy of thirty-five sequence mixer variants across six functional categories. Dense baselines provide full global routing at quadratic cost. Recurrent architectures enable constant-time inference through state compression. Sparse variants reduce complexity through structured patterns or token selection. Gated mechanisms introduce data-dependent control over memory retention and forgetting. Latent methods compress key--value representations into low-rank bottlenecks. Conditional memory provides test-time read/write over an external memory bank.}
\label{fig:comprehensive-taxonomy}
\end{figure}

\subsection{Dense and Recurrent Families}
The dense baselines (standard softmax \citep{vaswani_attention_2017} and sigmoid \citep{ramapuram_theory_2024}) and the recurrent/retentive family (RetNet \citep{sun_retentive_2023}, Mamba \citep{gu_mamba_2023}, ODE-style continuous blocks \citep{zhang_continuous_2021}, and Titans test-time memory \citep{behrouz_titans_2025}) provide the foundational expressiveness-versus-efficiency tradeoff. Mathematical formulations, parallel and recurrent forms, algorithmic pseudocode, and an architectural comparison table are collected in Appendix~\ref{sec:annex-transformer}. The pattern-driven mixer dispatch algorithm (including the training-free enforcement policy for \texttt{fasa\_attn} and \texttt{sparge\_attn}) is likewise detailed there.

\subsection{Latent and Low-Rank Attention Family}
A complementary family compresses the per-token key--value state into a low-rank latent: Multi-Head Latent Attention (MLA), Group-Query Latent Attention (GQLA), Multi-Head Low-Rank Attention (MLRA), Tucker Attention, Interleaved Head Attention (IHA), Grouped-head laTenT Attention (GTA), and Multi-head Temporal Latent Attention (MTLA). These methods unify grouped-query, latent, factorised, and temporal-latent designs under a common low-rank bottleneck. Full formulations, pseudocode, and a comparison table are collected in Appendix~\ref{sec:annex-latent}.

\subsection{Sparse Attention Extensions}
The mixer registry includes nine sparse blocks: \texttt{sparse\_transformer\_attn}, \texttt{longformer\_attn}, \texttt{bigbird\_attn}, \texttt{sparsek\_attn}, \texttt{nsa\_attn}, \texttt{sparge\_attn}, \texttt{fasa\_attn}, \texttt{msa\_attn}, and \texttt{sparda\_attn} \citep{child_sparse_transformer_2019,beltagy_longformer_2020,zaheer_bigbird_2020,lou_sparsek_2024,yuan_nsa_2025,zhang_spargeattn_2025,wang_fasa_2026,lai_minimax_2026,fu_sparda_2026}. The implementation enforces an explicit execution policy for training-free sparse methods: \texttt{fasa\_attn} and \texttt{sparge\_attn} are eval/inference-only and raise runtime errors if used while the model is in training mode. Mathematical formulations, algorithmic pseudocode, per-block characteristics, and a comparison table are collected in Appendix~\ref{sec:annex-sparse}.

\begin{figure}[t]
\centering
\fitfigure{\begin{tikzpicture}[node distance=0.9cm and 0.9cm]
\node[box, fill=gray!10] (tokens) {Token sequence};
\node[box, fill=blue!8, below left=1.1cm and 2.5cm of tokens] (local) {Local / windowed\\Longformer};
\node[box, fill=yellow!10, below=1.1cm of tokens] (hybrid) {Hybrid sparse graph\\BigBird / Sparse Transformer};
\node[box, fill=orange!12, below right=1.1cm and 2.5cm of tokens] (select) {Selective pruning\\SparseK / NSA / FASA / SpargeAttn};
\draw[line] (tokens) -- (local);
\draw[line] (tokens) -- (hybrid);
\draw[line] (tokens) -- (select);
\node[group, fit=(local)(hybrid)(select), label=below:{\small Three sparse design strategies: locality, structured graph sparsity, and learned or predicted selection}] {};
\end{tikzpicture}}
\caption{Conceptual map of sparse attention design choices used in the codebase. Different methods reduce cost by restricting neighborhoods, constructing sparse graphs, or selecting only high-value tokens/blocks.}
\label{fig:sparse-map}
\end{figure}

\begin{figure}[t]
\centering
\fitfigure{\begin{tikzpicture}[node distance=1.0cm and 1.2cm]
\node[box, fill=red!8] (standard) {Standard\\$\text{softmax}(QK^\top/\sqrt{d})V$};
\node[box, fill=red!8, right=of standard] (sigmoid) {Sigmoid\\$\sigma(QK^\top)V$};
\node[box, fill=red!8, right=of sigmoid] (gqa) {GQA\\Shared $K,V$ heads};
\node[group, fit=(standard)(sigmoid)(gqa), label=above:{\small Dense Attention: full pairwise interactions, $\mathcal{O}(n^2)$ cost}] {};
\end{tikzpicture}}
\caption{Three dense attention baselines. Standard softmax and sigmoid attention compute full $QK^\top$ matrices. Grouped-Query Attention (GQA) shares key--value heads across query groups, reducing memory bandwidth while preserving full attention expressivity.}
\label{fig:dense-map}
\end{figure}

\begin{figure}[t]
\centering
\fitfigure{\begin{tikzpicture}[node distance=1.0cm and 1.0cm]
\node[box, fill=green!8] (retnet) {RetNet\\Dual parallel/recurrent};
\node[box, fill=green!8, right=of retnet] (mamba) {Mamba\\Selective SSM};
\node[box, fill=green!8, right=of mamba] (ode) {ODE\\Continuous depth};
\node[box, fill=green!8, right=of ode] (titans) {Titans\\Test-time memory};
\node[group, fit=(retnet)(mamba)(ode)(titans), label=above:{\small Recurrent/Retentive: stateful, $\mathcal{O}(1)$ per-step inference}] {};
\end{tikzpicture}}
\caption{Recurrent and retentive architecture family. RetNet provides dual parallel and recurrent forms. Mamba uses selective state-space parameters. ODE blocks model continuous-depth dynamics. Titans introduce test-time memory adaptation. All achieve constant-time per-step inference.}
\label{fig:recurrent-map}
\end{figure}

\begin{figure}[t]
\centering
\fitfigure{\begin{tikzpicture}[node distance=1.0cm and 0.8cm]
\node[box, fill=cyan!8] (kv) {$K,V \in \mathbb{R}^{n \times d}$};
\node[box, fill=cyan!12, below=of kv] (down) {Down-project\\$W_{\text{down}} \in \mathbb{R}^{d \times r}$};
\node[box, fill=cyan!8, below=of down] (latent) {Latent $K',V' \in \mathbb{R}^{n \times r}$};
\node[box, fill=cyan!12, below=of latent] (up) {Up-project\\$W_{\text{up}} \in \mathbb{R}^{r \times d}$};
\node[box, fill=cyan!8, below=of up] (attn) {Attention\\$\text{softmax}(QK'^\top)V'$};
\draw[line] (kv) -- (down);
\draw[line] (down) -- (latent);
\draw[line] (latent) -- (up);
\draw[line] (up) -- (attn);
\node[group, fit=(kv)(down)(latent)(up)(attn), label=right:{\small Compression factor: $r \ll d$}] {};
\end{tikzpicture}}
\caption{Latent and low-rank attention concept. Key--value pairs are compressed through a down-projection into a low-rank latent space ($r \ll d$), where attention is computed at reduced cost, then up-projected back. Variants differ in factorization structure (MLA, GQLA, MLRA, Tucker), head interleaving (IHA, GTA), temporal compression (MTLA), and cross-covariance formulations (CCA, CCGQA).}
\label{fig:latent-map}
\end{figure}

\subsection{Gated Attention Extensions}
The mixer registry includes eight gated blocks: \texttt{gla\_attn}, \texttt{deltanet\_attn}, \texttt{gated\_deltanet\_attn}, \texttt{gated\_deltanet2\_attn}, \texttt{hgrn2\_attn}, \texttt{fox\_attn}, \texttt{gated\_softmax\_attn}, and \texttt{kda\_attn} \citep{yang_gla_2023,yang_deltanet_2024,yang_gated_deltanet_2024,hatamizadeh_gated_deltanet2_2026,qin_hgrn2_2024,lin_forgetting_transformer_2025,qiu_gated_attention_2025,kimi_kda_2025}. The unifying idea is that gating controls \emph{what information survives}: some gates act on recurrent state updates (GLA, DeltaNet variants, HGRN2), while others modify the full attention path itself (FoX and Gated Softmax). KDA introduces kernelized dynamic attention with learnable gating over the kernel bandwidth. This makes gating especially useful when the model must trade off recall, recency, and bounded memory. Mathematical formulations, algorithmic pseudocode, a generic gating template, and a comparison table are collected in Appendix~\ref{sec:annex-gated}.

\begin{figure}[t]
\centering
\fitfigure{\begin{tikzpicture}[node distance=1.2cm and 1.0cm]
\node[box, fill=green!8] (state) {Previous state\\$S_{t-1}$};
\node[box, fill=red!8, right=of state] (gate) {Gate(s)\\$\alpha_t,\beta_t,G_t,f_t$};
\node[box, fill=blue!8, right=of gate] (write) {New key/value\\or SDPA output};
\node[box, fill=yellow!12, below=of gate] (next) {Updated memory / output\\$S_t$ or $O_t$};
\draw[line] (state) -- (gate);
\draw[line] (gate) -- (write);
\draw[line] (state) |- (next);
\draw[line] (write) -- (next);
\end{tikzpicture}}
\caption{Generic gating template. A gate can decay existing memory, regulate write strength, or modulate dense attention outputs, depending on the block family.}
\label{fig:gating-template}
\end{figure}